\documentclass{article}

% Gestion des polices en fonction de l'encodage (PDFLaTeX, XeLaTeX, LuaTeX ou autre) :
\usepackage{ifxetex}
\usepackage{mathtools}
\usepackage{amssymb}
\usepackage{ifluatex}
\ifxetex
  % XeTeX est utilisé :
  \usepackage{fontspec}
\else
  \ifluatex
    % LuaTeX est utilisé :
    \usepackage{fontspec}
  \else
    % pdfLaTeX ou autre moteur est utilisé :
    \usepackage[T1]{fontenc}
    \usepackage[utf8]{inputenc}
    \usepackage{lmodern}
  \fi
\fi

% Autres packages :
\usepackage[french]{babel}
\usepackage{mdframed}
\usepackage{enumitem}
\usepackage{hyperref}
\usepackage{amsmath}
\usepackage{amsfonts}
\usepackage{stmaryrd}
\usepackage[explicit]{titlesec}
\usepackage{framed}
\usepackage{amsthm}
\usepackage{xcolor}
\usepackage{tcolorbox}
\tcbuselibrary{breakable}
% Package permettant de générer des graphes :
\usepackage{tikz}
\usetikzlibrary{arrows}
\usetikzlibrary{arrows.meta}
\usepackage{pgfplots}
\pgfplotsset{compat=1.18} % Préciser la version de pgfplots afin qu'il soit compatible avec le type de cocument (ici "article").
%package pour images :
\usepackage{graphicx}
\graphicspath{ {./ressources/} }
\sloppy

% --- Informations du document ---
\title{Rapport d'étude sur TOM TODO ALED}
\author{LE BER Tom \& PEROTTINO Tony}
\date{\today}

\begin{document}

\maketitle

% --- Sommaire ---
\newpage
\tableofcontents
\newpage

% --- Introduction ---
\section*{Préambule}
\addcontentsline{toc}{section}{Préambule}
Ce travail a été réalisé dans le cadre de la Double Licence Mathématiques-Informatique (MIDL) de l'université de Toulouse.
Il correspond à la matière "Stage MIDL 2" consistant en un travail dirigé de recherche (TER) dans les domaines de l'informatique et/ou des mathématiques.
Ce rapport fait office de compte-rendu de l'entièreté de notre travail de recherche.

% --- Corps du document ---
\section{Introduction au contexte de recherche}
\subsection{}



\end{document}